% =====================================================================
%  Thème 3 — Relation fondamentale & formules d'addition — NIVEAU MOYEN
%  Exercices
% =====================================================================
\documentclass[11pt,a4paper]{article}
\usepackage[utf8]{inputenc}
\usepackage[T1]{fontenc}
\usepackage{lmodern}
\usepackage[french]{babel}
\usepackage{amsmath,amssymb,mathtools}
\usepackage{icomma}
\usepackage{xcolor}
\usepackage{colortbl}
\usepackage{tikz}
\usepackage[most]{tcolorbox}
\usepackage{enumitem}
\usepackage{array}
\usepackage[a4paper,margin=2.2cm]{geometry}
\usetikzlibrary{arrows.meta,calc}

\definecolor{primary}{RGB}{214,51,108}
\definecolor{primarydark}{RGB}{140,20,64}

\DeclareMathOperator{\tg}{tg}
\DeclareMathOperator{\cotg}{cotg}
\newcommand{\degr}{^{\circ}}
\newcommand{\Z}{\mathbb{Z}}

\newtcolorbox{consigne}{enhanced,breakable,colback=primary!5,colframe=primary,
  boxrule=0.8pt,arc=2pt,left=3mm,right=3mm,top=2mm,bottom=2mm}

\newcounter{exo}
\newcommand{\exo}[1]{\medskip\refstepcounter{exo}%
  \noindent{\large\bfseries\color{primarydark}Exercice \theexo}%
  \ \textcolor{primary}{\textbullet}\ \textbf{#1}\par\nopagebreak\smallskip}

\setlength{\parindent}{0pt}
\pagestyle{empty}

\begin{document}

\begin{center}
{\LARGE\bfseries\color{primary}Thème 3 — Relation fondamentale \& formules d'addition}\\[2pt]
{\color{primarydark}\textsc{Niveau Moyen — Exercices}}
\end{center}
\hrule height 1pt
\medskip

\begin{consigne}
\textbf{Objectif :} à partir d'un seul nombre trigonométrique et du quadrant, retrouver les autres ;
obtenir des valeurs exactes par addition ; calculer un angle double. \textbf{Sans calculatrice},
valeurs \emph{exactes}, dénominateurs rationalisés.
\end{consigne}

\exo{Un sinus (ou cosinus) connu $\to$ le reste}
Pour chaque cas, calcule les deux autres nombres trigonométriques ($\cos\alpha$, $\sin\alpha$ ou
$\tg\alpha$), avec le bon signe.
\begin{enumerate}[label=\textbf{(\alph*)},leftmargin=2em,itemsep=2pt]
\item $\sin\alpha=\dfrac{3}{5}$, \ $\alpha$ au quadrant II. \quad Calcule $\cos\alpha$ et $\tg\alpha$.
\item $\sin\alpha=-\dfrac{1}{5}$, \ $\alpha$ au quadrant IV. \quad Calcule $\cos\alpha$ et $\tg\alpha$.
\item $\cos\alpha=-\dfrac{2}{3}$, \ $\alpha$ au quadrant III. \quad Calcule $\sin\alpha$ et $\tg\alpha$.
\end{enumerate}

\exo{Partir de la tangente}
On sait que $\tg\alpha=-\dfrac{3}{2}$ et que $\alpha$ appartient au quadrant II.
En utilisant $1+\tg^2\alpha=\dfrac{1}{\cos^2\alpha}$, détermine $\cos\alpha$ puis $\sin\alpha$.

\exo{Valeurs exactes par une différence ou une somme}
Décompose chaque angle en somme/différence d'angles remarquables, puis calcule la valeur exacte.
\[
\textbf{(a)}\ \cos15\degr \qquad
\textbf{(b)}\ \sin75\degr \qquad
\textbf{(c)}\ \cos105\degr
\]

\exo{La tangente d'un angle non remarquable}
À l'aide de $\tg(a\pm b)=\dfrac{\tg a\pm\tg b}{1\mp\tg a\,\tg b}$, calcule (résultat sans racine au
dénominateur).
\[
\textbf{(a)}\ \tg15\degr=\tg(45\degr-30\degr) \qquad
\textbf{(b)}\ \tg105\degr=\tg(60\degr+45\degr)
\]

\exo{Passer à l'angle double}
\begin{enumerate}[label=\textbf{(\alph*)},leftmargin=2em,itemsep=2pt]
\item $\sin\alpha=\dfrac{1}{3}$, $\alpha$ au quadrant I. Calcule $\cos\alpha$, puis $\sin2\alpha$ et
$\cos2\alpha$.
\item $\cos\alpha=-\dfrac{3}{5}$, $\alpha$ au quadrant II. Calcule $\sin\alpha$, puis $\sin2\alpha$ et
$\cos2\alpha$.
\end{enumerate}

\end{document}
